\documentclass[../../../include/open-logic-section]{subfiles}

\begin{document}

\olfileid{sth}{choice}{countablechoice}
\olsection{گِنن جوگ سیٹاں لئی چُون}

چُون یا حُسنِ ترتیب دے مسلّمے ورتے بناں ای ہِیٹھلی گل ثابت کرنی سوکھی اے:

\begin{lem}[$\Zminus$ وچ]
ہر محدود سیٹ لئی اک چُون والا فنکشن اے۔
\end{lem}

\begin{proof}
$a = \{b_1, \ldots, b_n\}$ ہووے۔ سادگی لئی منّ لو کہ ہر $b_i
\neq \emptyset$۔ سو کجھ شےواں $c_1, \ldots, c_n$ نیں جیہناں لئی
$c_1 \in b_1, \ldots, c_n \in b_n$۔ ہُن \olref[z][pairs]{prop:pairsconsequences} مطابق سیٹ $\{\langle b_1,
c_1\rangle , \ldots, \langle b_n, c_n\rangle\}$ موجود اے، تے ایہ~$a$ لئی اک چُون والا فنکشن اے۔
\end{proof}

پر لامتناہی سیٹ ویکھدیاں ای گل گنجل دار ہو جاندی اے۔ مثلاً اُتے دتے نتیجے دی ایہ «گھٹ توں گھٹ» توسیع ویکھو:

\begin{defish}
\emph{گِنن جوگ سیٹاں لئی چُون۔} ہر \emph{گِنن جوگ} سیٹ لئی اک چُون والا فنکشن اے۔
\end{defish}

ایہ چُون دے مسلّمے دی اک خاص صورت اے۔ تے پتا لگدا اے کہ ایہ اصول کافی واری ورتیا گیا سی، پر ورتن والیاں نوں ایہدا صاف پتا نہیں سی۔ دو سوہنیاں مثالاں ویکھو۔\footnote{ایہ مثالاں \citet[\S9.4]{Potter2004} تے لوکا اِنکُرواتی دیاں نیں۔}

\begin{ex}
اک سہج خیال ایہ اے۔ کسے وی سیٹ $A$ لئی یا $\cardle{\omega}{A}$، یا $\cardeq{A}{n}$ کسے $n \in \omega$ لئی۔ ایہ ایس سہج خیال نوں آکھن دا اک طریقہ اے کہ ہر سیٹ یا محدود اے یا لامتناہی۔ کانتور تے کئی ہور ریاضی داناں نے ایہ دعویٰ ثبوت دتے بناں کیتا سی۔ احتیاط ورتدیاں اسیں ایہ گل \olref[cardinals][classing]{generalinfinitycharacter} وچ ثابت کیتی سی۔ پر اوس ثبوت وچ اسیں $\ZFC$ اندر کم کردے ساں، کیوں جے اسیں منّ رہے ساں کہ کسے وی سیٹ $A$ نوں حُسنِ ترتیب دتی جا سکدی اے، تے ایس لئی $\card{A}$ دی موجودگی دی ضمانت اے۔ یعنی اسیں چُون دا مسلّمہ صاف منّیا سی۔

اصل وچ \citet{Dedekind1888} نے ایس دعوے دا اپنا ثبوت ایویں دِتا:

\begin{thm}[$\Zminus + \text{گِنن جوگ سیٹاں لئی چُون}$ وچ]
کسے وی $A$ لئی یا $\cardle{\omega}{A}$ یا $\cardeq{A}{n}$ کسے $n \in \omega$ لئی۔
\end{thm}

\begin{proof}
منّ لو $\cardneq{A}{n}$ سارے $n \in \omega$ لئی۔ تاں خاص طور اُتے ہر $n < \omega$ لئی اک ذیلی سیٹ $A_n \subseteq A$ اے جیہدے این ٹھیک $\cardexpo{2}{n}$ رکن نیں۔ ایس لڑی $A_0, A_1, A_2,
\ldots$ نوں ورت کے اسیں ہر $n$ لئی ایویں تعریف کردے آں:
\[
	B_n = A_n \setminus \bigcup_{i < n} A_i.
\]
ہُن ایہ ویکھو:
\begin{align*}
	\card{\bigcup_{i < n}A_n}
	&\leq \card{A_0} + \card{A_1} + \ldots + \card{A_{n-1}}\\
	&=1 + 2 + \ldots + 2^{n-1}\\
	& = 2^n - 1\\
	& < 2^n = \card{A_n}
\end{align*}
ایس لئی ہر $B_n$ دا گھٹ توں گھٹ اک رکن، $c_n$، اے۔ نالے $B_n$ والے سیٹ جوڑی جوڑی کرکے وکھرے نیں، سو جے $c_n = c_m$ تاں $n = m$۔ پر ہر $c_n
\in A$۔ سو فنکشن $f(n) = c_n$ اک انجیکشن $\omega \to A$ اے۔
\end{proof}
\noindent
ڈیڈیکنڈ نے ایہ نہیں دسیا سی کہ اوہنے گِنن جوگ سیٹاں لئی چُون دا اصول ورتیا۔ پر کی \emph{تسی} اوہدی ورتوں ویکھی؟ فیر ویکھو۔ (سچ مُچ \emph{فیر ویکھو}۔)

ثبوت وچ گِنن جوگ سیٹاں لئی چُون دا اصول دو واری ورتیا گیا اے۔ اک واری اسیں اپنے سیٹاں دی لڑی $A_0$، $A_1$، $A_2$، \dots\@ لین لئی ورتیا۔ فیر رکن $c_n$ ہر~$B_n$ وچوں چُنن لئی ورتیا۔ نالے چُون دی ایس ورتوں نوں کڈھیا نہیں جا سکدا۔ \citet[p.~138]{Cohen1966} نے ثابت کیتا کہ جے ساڈے کول چُون دا کوئی وی روپ نہیں، تاں نتیجہ نہیں چلدا۔ یعنی $\ZF$ نال ایہ گل بے ٹکراؤ اے کہ کجھ سیٹاں دا~$\omega$ نال \emph{موازنہ نہیں ہو سکدا}۔
\end{ex}

\begin{ex}
\citeyear{Cantor1878} وچ کانتور نے آکھیا سی کہ گِنن جوگ سیٹاں دی گِنن جوگ گنتی دا اتحاد وی گِنن جوگ اے۔ اوہنے ثبوت نہیں دتا، شاید ایس کرکے کہ ثبوت اوہنوں صاف لگدا سی۔ اسیں احتیاط ورتدیاں ایس نتیجے دا ہور عام روپ \olref[card-arithmetic][simp]{kappaunionkappasize} وچ ثابت کیتا سی۔ پر ساڈے ثبوت نے چُون دا مسلّمہ صاف منّیا سی۔ تے گھٹ عام نتیجے دے ثبوت نوں وی گِنن جوگ سیٹاں لئی چُون دا اصول چاہیدا اے۔

\begin{thm}[$\Zminus + \text{گِنن جوگ سیٹاں لئی چُون}$ وچ]
جے $A_n$ گِنن جوگ اے ہر $n \in \omega$ لئی، تاں $\bigcup_{n <
\omega} A_n$ گِنن جوگ اے۔
\end{thm}

\begin{proof}
عمومیت گنوائے بناں منّ لو کہ ہر $A_n \neq \emptyset$۔ سو ہر $n \in \omega$ لئی اک !!a{surjection} $f_n \colon \omega
\to A_n$ اے۔ $f \colon \omega \times \omega \to \bigcup_{n <
\omega} A_n$ دی تعریف $f(m, n) = f_n(m)$ نال دیو۔ نتیجہ نکلدا اے، کیوں جے $\omega
\times \omega$ گِنن جوگ اے (\olref[sfr][siz][zigzag]{natsquaredenumerable}) تے $f$ اک !!a{surjection} اے۔
\end{proof}
\noindent
کی تسی گِنن جوگ سیٹاں لئی چُون دی ورتوں ویکھی؟ ایہ ساڈے فنکشناں دی لڑی $f_0$، $f_1$، $f_2$، \dots\footnote{چُون دی ایہو جیہی ورتوں \olref[card-arithmetic][simp]{kappaunionkappasize} وچ وی ہوئی سی، جدوں اسیں آکھیا سی «ہر $\beta \in \cardfont{a}$ لئی اک !!a{injection}
$f_\beta$ مقرر کرو»۔} چُنن لئی ورتیا اے۔ تے ایتھے وی، جے چُون دا کوئی اصول نہ ہووے تاں نتیجہ نہیں چلدا۔ خاص طور اُتے \citet{FefermanLevy1963} نے ثابت کیتا کہ $\ZF$ نال ایہ گل بے ٹکراؤ اے کہ گِنن جوگ سیٹاں دی گِنن جوگ گنتی دے اتحاد دی کارڈینیلٹی~$\beth_1$ ہووے۔ پر رسل نے ایہو گل بڑی ہاسے والی طرح آکھی اے:
\begin{quote}
  ایس دی مثال اوہ لکھپتی اے جیہڑا جدوں وی اک جوڑی بوٹ خریددا سی نال جراباں دی اک جوڑی خریددا سی، تے ہور کدی نہیں۔ دوویں خریدن دا اوہنوں اینا شوق سی کہ آخر اوہدے کول بوٹاں دیاں $\aleph_0$ جوڑیاں تے جراباں دیاں $\aleph_0$ جوڑیاں ہو گئیاں\dots\@ بوٹاں وچ اسیں سجّے تے کھبّے دا فرق کر سکدے آں۔ سو ہر جوڑی وچوں اک چُن سکدے آں، یعنی سارے سجّے بوٹ یا سارے کھبّے بوٹ۔ پر جراباں لئی چُون دا کوئی ایہو جیہا اصول آپ نہیں سُجھدا۔ تے جدوں تک اسیں ضرب والا مسلّمہ [یعنی عملی طور اُتے چُون دا مسلّمہ] نہ منّیئے، اسیں یقین نہیں کر سکدے کہ کوئی ایہو جیہی جماعت اے جیہدے وچ ہر جوڑی وچوں اک جراب اے۔
  \citep[p.~126]{Russell1919}
\end{quote}
مختصر گل ایہ اے کہ ہِیٹھلی گل ثابت کرن لئی چُون دا کوئی روپ چاہیدا اے۔ جے تہاڈے کول جراباں دیاں گِنن جوگ گنتی جوڑیاں نیں، تاں تہاڈے کول (بس) گِنن جوگ گنتی جراباں نیں۔ تے اصل وچ گِنن جوگ سیٹاں لئی چُون دے اصول (یا اوہدے برابر دے اصول) توں بناں گِنن جوگ سیٹاں دی گِنن جوگ گنتی دا اتحاد گِنن جوگ ہون توں رہ سکدا اے۔
\end{ex}

سبق ایہ اے کہ گِنن جوگ سیٹاں لئی چُون دا اصول وار وار ورتیا گیا، پر ورتن والیاں نوں ایہدا بہت پتا نہیں سی۔ فلسفیانہ سوال ایہ اے۔ اسیں گِنن جوگ سیٹاں لئی چُون دا \emph{جواز} کیویں دے سکدے آں؟

سہج جواز دین دی اک کوشش وچ اسیں اک ایہو جیہے کم دی گل کر سکدے آں جیہدے لامتناہی قدم ہون۔ منّ لو اسیں پہلی چُون $\nicefrac{1}{2}$ منٹ وچ کردے آں، دوجی چُون $\nicefrac{1}{4}$ منٹ وچ، \dots، اپنی $n$ ویں چُون $\nicefrac{1}{2^n}$ منٹ وچ، \dots\@ تاں $1$~منٹ اندر اسیں چُونواں دی اک $\omega$ لڑی بنا لئی ہووے گی، تے اک چُون والا فنکشن تعریف کر لِتا ہووے گا۔

پر ایہ خیالی تجربہ سانوں واقعی کی دس سکدا اے؟ پہلاں تاں ایہ «چُنن» دے خیال نوں کافی لفظی معنی وچ لین اُتے ٹکدا اے۔ دوجا ایہ ریاضی نوں مابعدالطبیعیاتی امکان نال بنّھدا لگدا اے۔

ہور ضروری گل ایہ اے کہ ایہ کوشش \emph{بس} گِنن جوگ سیٹاں دی چُون توں اگے، \emph{پورے} چُون دے مسلّمے دا کوئی جواز نہیں دیندی۔ جے سانوں \emph{ہر} سیٹ لئی چُون والا فنکشن چاہیدا اے، تاں سانوں «کسے وی ترتیبی عدد جِنّی لمبائی والا لامتناہی قدماں دا کم» کر سکنا پئے گا۔ صاف آکھیے تاں ایہ ہاسے والی گل اے۔

\end{document}
